\section{Analisi dei task}

Spostando il focus dalle singole funzionalità dell'interfaccia (come l'apertura della mappa) ai reali obiettivi di fruizione, il peso di ogni operazione varia sensibilmente in base alle necessità di ciascun profilo. La tabella seguente classifica i macro-task distinguendo tra:
\begin{itemize}[nosep, leftmargin=1.2em]
    \item[\textcolor{culturiaRed}{\textbullet}] \textbf{Necessario (N)} — indispensabile per il raggiungimento del goal primario;
    \item[\textcolor{culturiaRed}{\textbullet}] \textbf{Facoltativo (F)} — arricchisce l'esperienza senza esserne vincolante.
\end{itemize}

% ─── Tabella Task (Parte 1) ──────────────────────────────────────
\renewcommand{\arraystretch}{1.25}

\begin{table}[H]
\centering
\small
\resizebox{\textwidth}{!}{%
\begin{tabular}{
    >{\raggedright\arraybackslash}p{4.2cm}
    >{\raggedright\arraybackslash}p{2.8cm}
    c c c
    >{\raggedright\arraybackslash}p{6.0cm}
}
\toprule
\rowcolor{culturiaRed}
\textcolor{culturiaWhite}{\textbf{Macro-task}} &
\textcolor{culturiaWhite}{\textbf{Persona}} &
\textcolor{culturiaWhite}{\textbf{Tipo}} &
\textcolor{culturiaWhite}{\textbf{Freq.}} &
\textcolor{culturiaWhite}{\textbf{Import.}} &
\textcolor{culturiaWhite}{\textbf{Motivazione}} \\
\midrule

% ── TASK 1 ─────────────────────────────────────────────────────
\rowcolor{culturiaBeige}
\textbf{1.\ Pianificazione e gestione del percorso} 
  & Marco {\footnotesize(T.\ Esterno)}     & \textbf{N} & Alta  & Alta  & Fondamentale per individuare rapidamente le opere principali nel limitato tempo a disposizione.  \\
\rowcolor{culturiaBeige}
  & Elena {\footnotesize(T.\ Int.\ Assiduo)} & \textbf{F} & Bassa & Bassa & Conosce già approfonditamente la planimetria e la disposizione delle sale.            \\
\rowcolor{culturiaBeige}
  & Matteo {\footnotesize(T.\ Int.\ Nuovo)}  & \textbf{N} & Media & Media & Indispensabile per orientarsi ed evitare il senso di smarrimento iniziale.           \\
\rowcolor{culturiaBeige}
  & Giovanni {\footnotesize(Funzionario)}      & \textbf{N} & Alta  & Alta  & Verifica che i percorsi e le mappe siano intuitivi per garantire l'orientamento dei flussi turistici comunali. \\

\midrule

% ── TASK 2 ─────────────────────────────────────────────────────
\textbf{2.\ Comprensione del patrimonio e approfondimento informativo}
  & Marco {\footnotesize(T.\ Esterno)}     & \textbf{F} & Bassa & Bassa & Predilige informazioni rapide e coincise, evitando testi troppo densi.         \\
  & Elena {\footnotesize(T.\ Int.\ Assiduo)} & \textbf{N} & Alta  & Alta  & Rappresenta il nucleo della visita: cerca retroscena, dettagli nascosti e restauri. \\
  & Matteo {\footnotesize(T.\ Int.\ Nuovo)}  & \textbf{N} & Media & Alta  & Cruciale per abbattere i tecnicismi attraverso la mediazione dell'IA. \\
  & Giovanni {\footnotesize(Funzionario)}      & \textbf{N} & Media & Alta  & Monitora la qualità e l'integrità culturale delle informazioni erogate per tutelare l'immagine dell'Ente. \\

\bottomrule
\end{tabular}%
}
\vspace{0.15cm}
\caption{Analisi dei macro-task orientati agli obiettivi (Parte 1): \textbf{N}\,=\,Necessario, \textbf{F}\,=\,Facoltativo.}
\label{tab:task_analysis_1}
\end{table}

% ─── Tabella Task (Parte 2) ──────────────────────────────────────
\begin{table}[H]
\centering
\small
\resizebox{\textwidth}{!}{%
\begin{tabular}{
    >{\raggedright\arraybackslash}p{4.2cm}
    >{\raggedright\arraybackslash}p{2.8cm}
    c c c
    >{\raggedright\arraybackslash}p{6.0cm}
}
\toprule
\rowcolor{culturiaRed}
\textcolor{culturiaWhite}{\textbf{Macro-task}} &
\textcolor{culturiaWhite}{\textbf{Persona}} &
\textcolor{culturiaWhite}{\textbf{Tipo}} &
\textcolor{culturiaWhite}{\textbf{Freq.}} &
\textcolor{culturiaWhite}{\textbf{Import.}} &
\textcolor{culturiaWhite}{\textbf{Motivazione}} \\
\midrule

% ── TASK 3 ─────────────────────────────────────────────────────
\rowcolor{culturiaBeige}
\textbf{3.\ Trasformare la passività in esplorazione tematica gratificante}
  & Marco {\footnotesize(T.\ Esterno)}     & \textbf{F} & Bassa & Bassa & Tempo limitato per concedersi deviazioni o logiche di gamification. \\
\rowcolor{culturiaBeige}
  & Elena {\footnotesize(T.\ Int.\ Assiduo)} & \textbf{N} & Alta  & Media & Ottimo incentivo per trasformare percorsi già noti in esplorazioni tematiche.   \\
\rowcolor{culturiaBeige}
  & Matteo {\footnotesize(T.\ Int.\ Nuovo)}  & \textbf{N} & Alta  & Alta  & La struttura a tappe e il sistema di ricompense stimolano la curiosità e vincono l'apatia. \\
\rowcolor{culturiaBeige}
  & Giovanni {\footnotesize(Funzionario)}      & \textbf{F} & Bassa & Bassa & Valuta l'efficacia delle logiche di gamification come leva per incrementare il tempo di permanenza sul territorio. \\

\midrule

% ── TASK 4 ─────────────────────────────────────────────────────
\textbf{4.\ Gestione e authoring dei contenuti culturali} 
  & Marco {\footnotesize(T.\ Esterno)}     & \textbf{F} & --    & --    & Non è di sua competenza; fruisce i contenuti senza crearli.               \\
  & Elena {\footnotesize(T.\ Int.\ Assiduo)} & \textbf{F} & --    & --    & Utente consumer: esplora i contenuti ma non ne è l'autore.                \\
  & Matteo {\footnotesize(T.\ Int.\ Nuovo)}  & \textbf{F} & --    & --    & Non pertinente al suo ruolo di visitatore occasionale.                    \\
  & Giovanni {\footnotesize(Funzionario)}      & \textbf{N} & Alta  & Alta  & Strumento essenziale per l'indipendenza dell'Ente: gestisce i contenuti in autonomia abbattendo i costi di outsourcing. \\

\bottomrule
\end{tabular}%
}
\vspace{0.15cm}
\caption{Analisi dei macro-task orientati agli obiettivi (Parte 2).}
\label{tab:task_analysis_2}
\end{table}

\vspace{0.8cm}

\subsection{Decomposizione dei Task Principali (Sub-Tasks)}

Al fine di calare l'analisi a un livello più concreto ed empirico, come richiesto dall'adozione dello \textit{User-Centered Design}, i macro-task sono stati formalmente spacchettati in \textbf{Sotto-Task} (Sub-Tasks). Questi descrivono i veri e propri passi operativi e le interazioni che gli utenti effettueranno sull'interfaccia.

% --- TABELLA 2 (Pianificazione) ---
\begin{table}[H]
\centering
\footnotesize
\renewcommand{\arraystretch}{1.0}
\resizebox{\textwidth}{!}{%
\begin{tabular}{
    >{\raggedright\arraybackslash}p{4.2cm}
    |>{\raggedright\arraybackslash}p{2.8cm}
    |>{\centering\arraybackslash}p{1.2cm}
    |>{\centering\arraybackslash}p{1.2cm}
    |>{\centering\arraybackslash}p{1.2cm}
    |>{\raggedright\arraybackslash}p{6.0cm}
}
\toprule
\rowcolor{culturiaRed}
\textcolor{culturiaWhite}{\textbf{Sotto-task}} &
\textcolor{culturiaWhite}{\textbf{Persona}} &
\textcolor{culturiaWhite}{\textbf{Tipo}} &
\textcolor{culturiaWhite}{\textbf{Freq.}} &
\textcolor{culturiaWhite}{\textbf{Import.}} &
\textcolor{culturiaWhite}{\textbf{Motivazione}} \\
\midrule

\textbf{1.1 Scansione RFID} \newline {\footnotesize(Avvio contestuale)}
  & Marco & \textbf{N} & Alta & Alta & Avvio immediato del tour senza navigazione manuale. \\\hline
  & Elena & \textbf{F} & Bassa & Bassa & Funzionalità non prioritaria durante le visite abituali. \\\hline
  & Matteo & \textbf{N} & Alta & Alta & Punto di ingresso intuitivo per chi non conosce il sistema. \\\hline
  & Giovanni & \textbf{N} & Alta & Alta & Verifica il corretto funzionamento delle postazioni RFID nei punti strategici del comune. \\

\midrule

\rowcolor{culturiaBeige}
\textbf{1.2 Individuazione di percorsi e servizi accessibili} 
  & Marco & \textbf{N} & Alta & Alta & Fondamentale per ottimizzare il percorso nel poco tempo concesso. \\\hline
\rowcolor{culturiaBeige}
  & Elena & \textbf{F} & Bassa & Bassa & Conosce già i percorsi principali privi di barriere. \\\hline
\rowcolor{culturiaBeige}
  & Matteo & \textbf{N} & Media & Alta & Supporto cruciale all'orientamento tramite indicazioni AR. \\\hline
\rowcolor{culturiaBeige}
  & Giovanni & \textbf{N} & Alta & Alta & Pianificazione di itinerari inclusivi per rendere il patrimonio accessibile a tutti i cittadini. \\

\midrule

\textbf{1.3 Sincronizzazione tour in simultanea}
  & Marco & \textbf{N} & Media & Alta & Coerente coordinazione della visita con il proprio gruppo. \\\hline
  & Elena & \textbf{F} & Bassa & Bassa & Visita solitamente in forma individuale. \\\hline
  & Matteo & \textbf{F} & Media & Media & Utile nel caso di visite in compagnia di colleghi. \\\hline
  & Giovanni & \textbf{F} & Bassa & Media & Analisi dei flussi per ottimizzare la gestione di gruppi numerosi ed evitare assembramenti. \\

\bottomrule
\end{tabular}%
}
\caption{Analisi dei Sotto-task: Pianificazione e gestione del percorso.}
\label{tab:subtasks_1}
\end{table}

\vspace{0.3cm}

% --- TABELLA 3 (Comprensione) ---
\begin{table}[H]
\centering
\footnotesize
\renewcommand{\arraystretch}{1.0}
\resizebox{\textwidth}{!}{%
\begin{tabular}{
    >{\raggedright\arraybackslash}p{4.2cm}
    |>{\raggedright\arraybackslash}p{2.8cm}
    |>{\centering\arraybackslash}p{1.2cm}
    |>{\centering\arraybackslash}p{1.2cm}
    |>{\centering\arraybackslash}p{1.2cm}
    |>{\raggedright\arraybackslash}p{6.0cm}
}
\toprule
\rowcolor{culturiaRed}
\textcolor{culturiaWhite}{\textbf{Sotto-task}} &
\textcolor{culturiaWhite}{\textbf{Persona}} &
\textcolor{culturiaWhite}{\textbf{Tipo}} &
\textcolor{culturiaWhite}{\textbf{Freq.}} &
\textcolor{culturiaWhite}{\textbf{Import.}} &
\textcolor{culturiaWhite}{\textbf{Motivazione}} \\
\midrule

\rowcolor{culturiaBeige}
\textbf{2.1 Scansione visiva ad overlay (AR)}
  & Marco & \textbf{F} & Bassa & Bassa & Poco tempo per interagire con layer informativi complessi. \\\hline
\rowcolor{culturiaBeige}
  & Elena & \textbf{N} & Alta & Alta & Indispensabile per scoprire dettagli e restauri nascosti. \\\hline
\rowcolor{culturiaBeige}
  & Matteo & \textbf{N} & Media & Alta & Stimolo visivo immediato per la comprensione dell'opera. \\\hline
\rowcolor{culturiaBeige}
  & Giovanni & \textbf{N} & Media & Alta & Ispezione periodica per validare la corretta sovrapposizione dei contenuti AR con il sito reale. \\

\midrule

\textbf{2.2 Selezione "Tone of Voice" (IA)}
  & Marco & \textbf{N} & Media & Alta & Predilige stile Storytelling per brevità e chiarezza. \\\hline
  & Elena & \textbf{N} & Alta & Alta & Utilizza lo stile accademico per scopi di studio. \\\hline
  & Matteo & \textbf{N} & Alta & Alta & Fondamentale per abbattere tecnicismi e godersi la visita. \\\hline
  & Giovanni & \textbf{N} & Media & Alta & Supervisione della qualità e del tono istituzionale delle risposte generate dall'IA. \\

\midrule

\rowcolor{culturiaBeige}
\textbf{2.3 Interrogazione vocale dinamica}
  & Marco & \textbf{F} & Bassa & Bassa & Preferisce non fermarsi per query vocali specifiche. \\\hline
\rowcolor{culturiaBeige}
  & Elena & \textbf{N} & Alta & Alta & Domande verticali su dettagli specifici dell'opera. \\\hline
\rowcolor{culturiaBeige}
  & Matteo & \textbf{N} & Media & Media & Interazione naturale senza navigazione manuale dei menu. \\\hline
\rowcolor{culturiaBeige}
  & Giovanni & \textbf{F} & Bassa & Bassa & Valuta l'interazione vocale come strumento per abbattere le barriere informative del borgo. \\

\bottomrule
\end{tabular}%
}
\caption{Analisi dei Sotto-task: Comprensione del patrimonio e approfondimento informativo.}
\label{tab:subtasks_2}
\end{table}

\newpage

% --- TABELLA 4 (Gamification) ---
\begin{table}[H]
\centering
\small
\renewcommand{\arraystretch}{1.3}
\resizebox{\textwidth}{!}{%
\begin{tabular}{
    >{\raggedright\arraybackslash}p{4.2cm}
    |>{\raggedright\arraybackslash}p{2.8cm}
    |>{\centering\arraybackslash}p{1.2cm}
    |>{\centering\arraybackslash}p{1.2cm}
    |>{\centering\arraybackslash}p{1.2cm}
    |>{\raggedright\arraybackslash}p{6.0cm}
}
\toprule
\rowcolor{culturiaRed}
\textcolor{culturiaWhite}{\textbf{Sotto-task}} &
\textcolor{culturiaWhite}{\textbf{Persona}} &
\textcolor{culturiaWhite}{\textbf{Tipo}} &
\textcolor{culturiaWhite}{\textbf{Freq.}} &
\textcolor{culturiaWhite}{\textbf{Import.}} &
\textcolor{culturiaWhite}{\textbf{Motivazione}} \\
\midrule

\textbf{3.1 Dashboard e Progressioni}
  & Marco & \textbf{F} & Bassa & Bassa & Nessun interesse per progressioni di lungo periodo. \\\hline
  & Elena & \textbf{F} & Media & Bassa & Tracciamento facoltativo delle tappe completate. \\\hline
  & Matteo & \textbf{N} & Alta & Alta & Motore per sconfiggere l'inerzia e l'apatia iniziale. \\\hline
  & Giovanni & \textbf{F} & Bassa & Bassa & Analisi dei dati di completamento per monitorare l'interesse verso i percorsi tematici. \\

\midrule

\rowcolor{culturiaBeige}
\textbf{3.2 Rewarding tramite tap RFID}
  & Marco & \textbf{F} & Bassa & Bassa & Poca attrattiva per i punti data la visita singola. \\\hline
\rowcolor{culturiaBeige}
  & Elena & \textbf{F} & Media & Bassa & Gradisce il riconoscimento ma non ne è dipendente. \\\hline
\rowcolor{culturiaBeige}
  & Matteo & \textbf{N} & Alta & Alta & Feedback immediato che rinforza l'esplorazione attiva. \\\hline
\rowcolor{culturiaBeige}
  & Giovanni & \textbf{F} & Bassa & Bassa & Valuta l'impatto dei premi digitali sul coinvolgimento della comunità locale. \\

\midrule

\textbf{3.3 Liste di Missioni contestuali}
  & Marco & \textbf{F} & Bassa & Bassa & Missioni troppo articolate per lo scarso tempo disponibile. \\\hline
  & Elena & \textbf{N} & Media & Alta & Incentivo a esplorare zone del polo solitamente ignorate. \\\hline
  & Matteo & \textbf{N} & Alta & Alta & Guida a tappe che struttura l'intera visita al museo. \\\hline
  & Giovanni & \textbf{F} & Alta & Bassa & Strumento per incentivare la scoperta di aree del borgo solitamente meno visitate. \\

\bottomrule
\end{tabular}%
}
\caption{Analisi dei Sotto-task: Trasformare la passività in esplorazione gratificante.}
\label{tab:subtasks_3}
\end{table}

\vspace{0.3cm}

% --- TABELLA 5 (Authoring / Amministratore) ---
\begin{table}[H]
\centering
\footnotesize
\renewcommand{\arraystretch}{1.0}
\resizebox{\textwidth}{!}{%
\begin{tabular}{
    >{\raggedright\arraybackslash}p{4.2cm}
    |>{\raggedright\arraybackslash}p{2.8cm}
    |>{\centering\arraybackslash}p{1.2cm}
    |>{\centering\arraybackslash}p{1.2cm}
    |>{\centering\arraybackslash}p{1.2cm}
    |>{\raggedright\arraybackslash}p{6.0cm}
}
\toprule
\rowcolor{culturiaRed}
\textcolor{culturiaWhite}{\textbf{Sotto-task}} &
\textcolor{culturiaWhite}{\textbf{Persona}} &
\textcolor{culturiaWhite}{\textbf{Tipo}} &
\textcolor{culturiaWhite}{\textbf{Freq.}} &
\textcolor{culturiaWhite}{\textbf{Import.}} &
\textcolor{culturiaWhite}{\textbf{Motivazione}} \\
\midrule

\textbf{4.1 Creazione/modifica scheda opera}
  & Giovanni & \textbf{N} & Alta  & Alta  & Inserimento e aggiornamento di testi, immagini e tag AR per ogni opera esposta. \\\hline

\rowcolor{culturiaBeige}
\textbf{4.2 Configurazione percorso AR guidato}
  & Giovanni & \textbf{N} & Media & Alta  & Definizione dell'ordine delle tappe, aree trigger RFID e contenuto AR associato a ciascuna. \\\hline

\textbf{4.3 Monitoraggio statistiche di visita}
  & Giovanni & \textbf{N} & Alta & Alta & Analisi dei dati d'uso per giustificare futuri investimenti o variazioni nell'offerta culturale. \\

\bottomrule
\end{tabular}%
}
\caption{Analisi dei Sotto-task: Gestione e authoring dei contenuti culturali (Giovanni).}
\label{tab:subtasks_4}
\end{table}

\vspace{0.6cm}

\textbf{Nota sullo sviluppo iterativo:}\\
L'elenco dei task qui riportato rappresenta una base progettuale preliminare. I questionari di analisi (rivolti sia ai turisti che agli addetti al settore) sono tuttora in fase di acquisizione: i feedback e le idee che emergeranno da questa fase di ascolto attivo degli stakeholder porteranno a integrazioni e modifiche dirette nel prossimo \textbf{Assignment 2}, seguendo la metodologia iterativa tipica dello \textit{User-Centered Design}.
